\documentclass[tikz,border=10pt]{standalone}
\usepackage[utf8]{inputenc}
\usepackage[T1]{fontenc}
\usepackage[spanish]{babel}
\usepackage{tikz}
\usetikzlibrary{positioning,arrows.meta}

\definecolor{cafeborder}{HTML}{4A2F24}

\tikzset{
  box/.style={
    draw=cafeborder,
    line width=0.8pt,
    rounded corners=3pt,
    align=center,
    minimum height=10mm,
    text width=38mm,
    inner sep=5pt,
    font=\sffamily
  },
  rootbox/.style={
    box,
    line width=1.2pt,
    text width=48mm,
    minimum height=12mm,
    font=\sffamily\bfseries
  },
  groupbox/.style={
    box,
    text width=40mm,
    font=\sffamily\bfseries
  },
  statebox/.style={
    box,
    text width=36mm
  },
  relation/.style={
    -{Stealth[length=2.2mm]},
    draw=cafeborder,
    line width=0.8pt
  }
}

\begin{document}
\begin{tikzpicture}[node distance=12mm and 16mm]

  % Nodo raíz y área funcional
  \node[rootbox] (app) {CaféU\\Aplicación móvil};
  \node[groupbox, below=of app] (order) {Pedido};

  % Pantallas principales
  \node[box, below left=of order, xshift=-22mm] (menu) {Menú};
  \node[box, below=of order] (review) {Revisión del pedido};
  \node[groupbox, below right=of order, xshift=22mm] (status) {Estado del pedido};

  % Estados del envío
  \node[statebox, below left=of status, xshift=-12mm] (sending) {Enviando pedido\\\footnotesize Estado transitorio};
  \node[statebox, below=of status] (received) {Pedido recibido\\\footnotesize Resultado exitoso};
  \node[statebox, below right=of status, xshift=12mm] (error) {Error de envío\\\footnotesize Estado recuperable};

  % Relaciones estructurales: las flechas indican agrupación, no secuencia
  \draw[relation] (app) -- (order);
  \draw[relation] (order) -- (menu);
  \draw[relation] (order) -- (review);
  \draw[relation] (order) -- (status);
  \draw[relation] (status) -- (sending);
  \draw[relation] (status) -- (received);
  \draw[relation] (status) -- (error);

  % Nota de lectura
  \node[below=17mm of review, text width=125mm, align=center, font=\sffamily\small]
    (note) {Las flechas muestran la organización de pantallas y estados; no representan el recorrido paso a paso del usuario.};

\end{tikzpicture}
\end{document}

